\begin{table*}[t]
\centering
\caption{Composition and provenance of ProactiveStream. Reference interventions mark service opportunities.}
\label{tab:benchmark_stats}
\resizebox{\textwidth}{!}{
\begin{tabular}{@{}lcccccc@{}}
\toprule
scope  & primary source & entity          & streams & \begin{tabular}[c]{@{}l@{}}input\\ events\end{tabular} & temporal coverage                 & \begin{tabular}[c]{@{}l@{}}reference\\ interventions\end{tabular} \\ \midrule
Sports & HeartSteps     & participant     & 9       & 1,456                                                  & 367 observed participant-days     & 365                                                               \\
Home   & CASAS          & home            & 4       & 1,624                                                  & 14 consecutive days/home          & 190                                                               \\
Code   & GH Archive     & repository      & 4       & 1,459                                                  & longitudinal repository histories & 174                                                               \\\midrule
total  & three scopes   & isolated stream & 17      & 4,539                                                  & three longitudinal settings       & 729                                                               \\ \bottomrule
\end{tabular}
}
\end{table*}

\section{ProactiveStream benchmark}
\label{sec:benchmark}

ProactiveStream evaluates service-opportunity detection in longitudinal streams as agents accumulate history.
Table~\ref{tab:benchmark_stats} gives its composition.

\subsection{Scopes and public events}

\hspace*{\parindent}\textbf{Sports} uses participant activity and availability contexts for optional walking and everyday physical-activity assistance.

\textbf{Home} uses time-bounded home-sensor observations for gentle, low-risk check-ins.
Each sensor segment becomes available at interval closure.

\textbf{Code} uses repository activity at GitHub's native event granularity for low-risk triage, review, and maintenance, with one decision per event.

Each participant, home, or repository defines an isolated chronological stream, allowing earlier observations to inform later service opportunities.
Public inputs contain contemporaneously available observations, using the same scope-specific schema for both label classes.
Agents construct their own memory from these events; annotation and scoring metadata remain evaluator-private.
Appendix~\ref{app:benchmark_protocol} details source datasets, fields, and preprocessing.

\subsection{Reference labels and causal replay}
\label{sec:reference_labels}

Annotators assess $e_t$ and its complete prior same-stream history $\mathcal{H}_{<t}$.
Let $A$ indicate a concrete, timely, feasible, low-risk opportunity supported by $e_t$, $B$ an opportunity at $e_t$ supported by historical observations, and $N_t^{\mathrm{ref}}$ reference-episode novelty shared across methods:
\begin{equation}
\begin{aligned}
y_t=1 \quad \Longleftrightarrow \quad &
\big(A(e_t)=1\ \lor\ B(e_t,\mathcal{H}_{<t})=1\big)\\
&\land N_t^{\mathrm{ref}}=1.
\end{aligned}
\label{eq:or_semantics}
\end{equation}
The reference protocol labels the first reliable actionable window \emph{intervene} and covered windows \emph{abstain}, allowing renewed opportunities when support strengthens, the service changes, or an episode reopens.
Event-level labels assess timely detection of these service windows.

LLM-assisted triage prioritizes candidate opportunities; three team members then independently label every event while blinded to evaluated-system outputs, followed by cross-review and adjudication.
Appendix~\ref{app:label_review_450} details a separate-team diagnostic of reference sensitivity and episode behavior.
Main comparisons fix labels, prompts, memory bounds, and validation settings before evaluation.

Each stream is replayed chronologically from empty method state, with access to the current event and strictly prior same-stream context.
Each decision is finalized before the current event is committed to memory; reference labels are joined after inference for scoring.
Scoring includes the initial events and continues as history accumulates.
